% SPDX-License-Identifier: Apache-2.0
\documentclass[11pt]{article}

\usepackage[margin=1in]{geometry}
\usepackage{amsmath,amssymb}
\usepackage{booktabs}
\usepackage{microtype}
\usepackage[hidelinks]{hyperref}
\usepackage{xcolor}

\newcommand{\sat}{\operatorname{sat}_{16}}
\newcommand{\ceilleak}[1]{\left\lceil |#1| / 2^{k} \right\rceil}

\title{\textbf{Mycelium: A Verification-First Spiking Network\\from the CeliumNeUR Constraint Set}}
\author{Mario Guti\'errez\\Celiums Solutions LLC\\\texttt{mario@celiums.ai}}
\date{August 2026 --- working draft v0.2}

\begin{document}
\maketitle

\begin{abstract}
Most spiking neural network (SNN) papers propose neuron dynamics and assert
hardware efficiency conditionally, against silicon that does not exist. We
invert the direction: we take CeliumNeUR --- a verification-first neuromorphic
SoC with synthesizable RTL, a bit-exact golden model, mutation testing and
bounded formal checks --- and derive a trainable model \emph{from its
constraint set}: ceiling leak, saturating 16-bit membranes, ternary saturating
weights, per-neuron thresholds, leak and refractory periods, tick-synchronous
phase semantics, and static routing. The result, \textbf{Mycelium}, is a
recurrent block-sparse ternary architecture whose forward pass is bit-identical
to the chip's verified golden model under an explicit five-condition contract,
whose learned topology is frozen into a hashable, input-independent artifact,
and whose deployment instance --- 34 neurons and 832 synapse entries, inside
the v1 silicon budget --- replays through the integer reference with zero
mismatches after training on real data. On Spiking Heidelberg Digits, Mycelium
beats a budget-matched fp16 GRU on the memory frontier in the edge regime
($62.3\% \pm 3.0$ at 31\,KB versus $55.4\% \pm 4.0$ at 68\,KB, three seeds
each), where the per-operation frontier also tilts in its favor; dense wins
above ${\sim}300$\,KB and in raw operation count elsewhere --- exactly the
split a memory-bandwidth thesis predicts. Along the way, $L_0$ topology
learning \emph{deleted the recurrent synapse entirely} when the task did not
pay for it; a float-weight ablation \emph{lost} to ternary under identical
dynamics; and a membrane-potential readout, chip-honest through the RTL's
non-invasive readback port, improved both accuracy ($+2.6$ points) and seed
stability ($\pm 1.9$ vs.\ $\pm 2.9$) over a spike-count head.
\end{abstract}

\section{Introduction: constraints first, hardware later}

BitNet's ternary weights were motivated by multiplication-free hardware that
GPUs lack; BitNet wins anyway --- on memory bandwidth
\cite{ma2024bitnet}. We claim the same structure holds for neuromorphic
constraint sets. Split the chip's efficiency story in two:

\begin{itemize}
  \item a \textbf{static half} --- sparse frozen topology and integer weights
    --- which is a \emph{memory} property and transfers to commodity
    accelerators today;
  \item a \textbf{dynamic half} --- event-driven ``no spike, no work''
    execution --- which is a \emph{compute} property, worthless below a block
    granularity of ${\sim}64$ on GPUs, and latent until silicon exists.
\end{itemize}

Designing against the full constraint set from the first training step
matters: a model trained with relaxed float dynamics (exponential leak, long
subthreshold memory) and quantized afterward depends on memory the silicon
does not have, and fails on deployment.

A second claim is unique to this project: \textbf{verifiability as a
feature}. A frozen compute graph is an artifact --- hashable, signable,
publishable. Because CeliumNeUR ships a golden model that its RTL is verified
against cycle-by-cycle (mutation testing: 17/17 killed; bounded model
checking to depth 60; lint-clean synthesis receipts), a network that is
bit-exact against that golden model inherits a chain of custody no SNN paper
we know of can offer:
\[
\text{task-trained weights} \;\equiv\; \text{integer model} \;\equiv\;
\text{golden referee} \;\equiv\; \text{verified RTL.}
\]

\section{The constraint set}

CeliumNeUR v1 is a 4-tile, 1{,}024-neuron, 1{,}024-synapse-entry SoC
(2$\times$2 mesh, X-first dimension-order routing, credit-based flow
control). The authority order for semantics is: golden model $>$
specification prose. The neuron, extracted from the golden source:

\paragraph{State.} Membrane $v \in [-32768, 32767]$ (signed 16-bit,
\emph{saturating}, never wrapping); refractory countdown $c \geq 0$.
Per-neuron parameters (all learnable in Mycelium, all stored per neuron on
chip): threshold $\theta \in [1, 32767]$, leak shift $k \in [0,15]$,
refractory duration $R$, reset mode.

\paragraph{One global tick.} With phase input sum $I_t$ (spikes staged at
tick $t{-}1$, integer counts $\times$ integer weights) and event-presence
mask $E_t$:
\begin{align}
  v' &= \sat(v_{t-1} + I_t) \\
  f^{\mathrm{evt}}_t &= E_t \wedge (c_{t-1} = 0) \wedge (v' \geq \theta) \\
  v'' &= v' - f^{\mathrm{evt}}_t\,\theta, \qquad
        c' = \begin{cases} R & f^{\mathrm{evt}}_t \\ c_{t-1} & \text{else} \end{cases} \\
  v''' &= \sat\!\big(v'' - \operatorname{sign}(v'')\,\ceilleak{v''}\big) \\
  f^{\mathrm{tick}}_t &= (c' = 0) \wedge (v''' \geq \theta) \\
  v_t &= v''' - f^{\mathrm{tick}}_t\,\theta, \qquad
        c_t = \max\!\big(c'[f^{\mathrm{tick}}_t \mapsto R] - 1,\; 0\big)
\end{align}
The ceiling leak loses at least one unit per tick: subthreshold memory is
\emph{proportional to signal magnitude}, never exponential-tailed. The
event gate $E_t$ is load-bearing (a zero-weight synapse still evaluates on
chip; a quiet phase evaluates only post-leak), and the
decrement-after-evaluation order makes the two fire paths deliberately
asymmetric (an event-path fire blocks $R$ ticks; a tick-path fire blocks
$R{-}1$) --- both discovered during extraction and since documented upstream
as contract.

\paragraph{The equivalence contract (C1--C5).} The golden model evaluates
fire per synaptic event; a batched model evaluates per tick. Exact
end-of-phase state equality holds under: \textbf{C1} subtractive reset,
\textbf{C2} $R \geq 1$, \textbf{C3} no intermediate 16-bit saturation,
\textbf{C4} no order-dependent threshold crossing, \textbf{C5} external
injection subthreshold at injection time. Every boundary is pinned by an
explicit divergence test. Under the contract we verified exact spike and
state equality on over $10^5$ fuzzed neuron-phases and on full 1{,}024-neuron
chip simulations (mesh, dendrite tables with true duplicate multiplicity,
somas), with delivery routed through the same sparse primitive used for
training.

\section{Mycelium}

\paragraph{Macro-architecture.} One or more recurrent modules over a
\emph{static block graph}: feed and recurrent synapses are block-sparse
ternary matrices (an active $B{\times}B$ block is all $B^2$ dendrite entries,
making event semantics exact and the mask a faithful FLOP accountant),
neurons are the integer dynamics above with a straight-through/surrogate
training path, and a linear head reads the top layer's \emph{membrane
potential in threshold units} $v/\theta$ at every tick, averaged. The
membrane head is chip-honest: the RTL exposes soma state through a
non-invasive readback port (invariant I5), so the head is a host-side layer;
a spike-count head is the fully-on-chip alternative at $-2.6$ points and
$1.5\times$ the seed variance. No attention, no data-dependent routing: the
compute graph is input-independent and hashable.

\paragraph{Training.} Surrogate gradients (arctan kernel, width
$0.25\,\theta$) at both comparators; straight-through estimators through the
ceiling leak, weight round\-ing, and the threshold grid; hard saturation
passes no gradient beyond the rails. Weights are ternarized per block by
absmean \cite{ma2024bitnet}; thresholds are float masters STE-rounded to the
chip grid; leak is heterogeneous per neuron across the full 4-bit range
\cite{perez2021neural}. Topology is learned with hard-concrete $L_0$ gates
\cite{louizos2018learning} over the block grid at a hot learning rate, then
\emph{binarized and frozen}; fine-tuning re-accommodates weights to the fixed
graph, whose masks hash to a SHA-256 artifact. Adam with cosine decay; spike
dropout; no gradient clipping (it hurt at every strength we tried).

\section{Results}

All numbers reproducible from the repository's \texttt{experiments/} JSONs.

\subsection{Gradient viability (pre-registered gate)}

A 324-configuration sweep (surrogate shape $\times$ width $\times$ weight
precision $\times$ $T$ $\times$ leak $\times$ threshold placement) passed
every pre-registered kill criterion. A previously observed null result
(${\sim}2\times10^{-5}$ gate gradient in a ternary predecessor) reproduced
\emph{only} in the sweep's worst corner --- fast leak, high threshold,
compact-support (triangular) kernel --- identifying it as a
threshold-placement pathology rather than a ternarization cost. Ternary was
the \emph{best}-conditioned precision (median relative input-layer gradient
$3.1\times10^{-2}$, vs.\ float $6.8\times10^{-3}$ and int8
$2.4\times10^{-4}$); fat-tailed kernels are robust where triangular dies;
slow leak dominates trainability; BPTT carries credit at $T{=}32$ without
vanishing.

\subsection{Topology learning}

On a task with a planted-relevance structure (one signal input block, one
pure-noise block), $L_0$ gating closed every noise route exactly at zero
accuracy cost, and random masks at the learned density collapsed to chance
half the time. On SHD, the learner made an unprompted structural decision:
at $\lambda = 0.15$ it \textbf{deleted the recurrent synapse entirely}
($-1.9$ points, $3\times$ fewer operations) --- rate-readout SHD at $T{=}32$
does not pay for recurrence, and the method discovered that by itself.

\subsection{Task curves}

Spiking Heidelberg Digits \cite{cramer2020heidelberg} (20 classes, 700 spike
channels, integer count binning at $T{=}32$), against fp16 GRUs trained under
the same budget (60 epochs cosine, learning-rate swept). \textbf{Every row:
three seeds.}

\begin{table}[h]
\centering
\begin{tabular}{lccc}
\toprule
Model & Accuracy & Bytes & MMAC/sample \\
\midrule
\textbf{Mycelium $\lambda{=}0.15$} & $\mathbf{0.623 \pm 0.030}$ & \textbf{31\,KB} & \textbf{1.5} \\
GRU-16 (fp16) & $0.554 \pm 0.040$ & 68\,KB & 1.1 \\
GRU-32 (fp16) & $0.669 \pm 0.033$ & 140\,KB & 2.3 \\
Mycelium $\lambda{=}0.02$ & $0.673 \pm 0.019$ & 167\,KB & 19.3 \\
Mycelium $H{=}1024$ & $0.650 \pm 0.021$ & 284\,KB & 32.1 \\
GRU-256 (fp16, $T{=}64$+aug) & $0.883 \pm 0.007$ & 1453\,KB & 47.5 \\
\bottomrule
\end{tabular}
\caption{SHD frontier. Bytes: 2-bit ternary weights in active blocks + block
masks + per-neuron soma parameters vs.\ fp16 dense. MACs count ternary
additions as multiplies and ignore ${\sim}0.4$ hidden event sparsity --- both
conservative against Mycelium.}
\end{table}

\textbf{The byte frontier favors Mycelium below ${\sim}150$\,KB} ($+6.9$
points at less than half the bytes at the extreme edge; a tie near 150\,KB)
and dense above ${\sim}300$\,KB. \textbf{The per-operation frontier crosses
in the extreme-edge corner}: at 1.5\,MMAC Mycelium scores $0.623$ where the
dense interpolation between its bracketing points passes ${\approx}0.59$;
elsewhere dense wins per-op. Doubling temporal resolution and adding
augmentation moved the GRU $+5$ points and Mycelium none: the remaining
accuracy gap at scale is architectural, of which the membrane head recovered
$+2.6$ points while halving seed variance. Multi-seeding also showed the
small dense baselines to be equally noisy (GRU-16: $\pm 4.0$): at this scale,
seed sensitivity is a property of the regime, not of spiking. DVS128 Gesture
\cite{amir2017dvs} (11 classes, chance 9.1\%) replicates the pattern with the
spike head: $0.689$ at 133\,KB vs.\ GRU $0.773$ at 484\,KB.

\subsection{The deployment certificate}

A two-class SHD instance inside the v1 silicon budget --- 34 neurons and 832
dendrite entries, including the host-side encoder --- trains to $74.9\%$ and
replays through the integer model (bit-exact to the golden referee, hence to
RTL semantics) with \textbf{zero logit mismatches} over the full test set.
The artifact is a SHA-256 over edges, int8 weights and integer thresholds.
To our knowledge this is the first task-trained SNN with a machine-checkable
equivalence chain to silicon-verified RTL.

\section{Related work}

ODIN \cite{frenkel2019odin} and ReckOn \cite{frenkel2022reckon} define the
digital-neuromorphic baseline CeliumNeUR's own specification calibrates
against. snnTorch \cite{eshraghian2023training} supplies the
surrogate-gradient canon; we deliberately did not fork it --- a tested
adapter exposes our neuron through its interface instead. Spikformer
\cite{zhou2023spikformer} shows spiking attention exists, but attention is
data-dependent routing --- precisely what a signable static graph excludes.
BitNet b1.58 \cite{ma2024bitnet} is the structural precedent for the memory
thesis; hard-concrete $L_0$ \cite{louizos2018learning} provides the gate
machinery; neural heterogeneity \cite{perez2021neural} motivates per-neuron
time constants, which the chip stores natively. Cramer et
al.~\cite{cramer2020heidelberg} report recurrent-LIF baselines around
$0.71$ at higher $T$, situating our $0.62$--$0.67$ at $T{=}32$ as in-family
for spiking models, not state of the art. Surrogate-gradient learning
follows \cite{neftci2019surrogate,zenke2021remarkable}; the arctan kernel is
from \cite{fang2021incorporating}.

\section{Limitations}

Inference-side efficiency only (BPTT training costs more than dense); byte
accounting at 2 bits/weight (1.58 information-theoretic would shift the
Mycelium curve further left); one and a half tasks, with the membrane head
unevaluated on DVS; three seeds per point (honest error bars, not
significance theater --- the edge-point gap is ${\sim}2\sigma$); no GPU
block-sparse kernel wall-clock measurements (FLOP fractions stand in); the
chip-faithful instance is a deployability certificate, not a competitive
model; scaling past 1{,}024 neurons on-silicon requires a flit redesign.

\section{Conclusion}

Constraint-first design paid where the thesis said it would: memory. The
verification chain --- model to silicon referee, bit-exact, signed --- is the
claim we recommend leading with; the edge-regime frontier win is the number
that makes it more than a curiosity.

\paragraph{Artifacts.} Model: \texttt{github.com/terrizoaguimor/celiumsnn}
(Apache-2.0). Chip: \texttt{github.com/terrizoaguimor/celiumneur}
(Apache-2.0), DOI \texttt{10.5281/zenodo.21925426}.

\begin{thebibliography}{12}
\bibitem{cramer2020heidelberg} B. Cramer, Y. Stradmann, J. Schemmel,
F. Zenke. The Heidelberg spiking data sets for the systematic evaluation of
spiking neural networks. \emph{IEEE TNNLS}, 2020.
\bibitem{frenkel2019odin} C. Frenkel, M. Lefebvre, J.-D. Legat, D. Bol.
A 0.086-mm$^2$ 12.7-pJ/SOP 64k-synapse 256-neuron online-learning digital
spiking neuromorphic processor in 28-nm CMOS. \emph{IEEE TBioCAS}, 13(1), 2019.
\bibitem{frenkel2022reckon} C. Frenkel, G. Indiveri. ReckOn: A 28nm
sub-mm$^2$ task-agnostic spiking recurrent neural network processor enabling
on-chip learning over second-long timescales. \emph{ISSCC}, 2022.
\bibitem{eshraghian2023training} J. K. Eshraghian et al. Training spiking
neural networks using lessons from deep learning. \emph{Proc.\ IEEE},
111(9), 2023.
\bibitem{zhou2023spikformer} Z. Zhou et al. Spikformer: When spiking neural
network meets transformer. \emph{ICLR}, 2023.
\bibitem{ma2024bitnet} S. Ma et al. The era of 1-bit LLMs: All large
language models are in 1.58 bits. arXiv:2402.17764, 2024.
\bibitem{louizos2018learning} C. Louizos, M. Welling, D. P. Kingma.
Learning sparse neural networks through $L_0$ regularization. \emph{ICLR},
2018.
\bibitem{perez2021neural} N. Perez-Nieves, V. C. H. Leung, P. L. Dragotti,
D. F. M. Goodman. Neural heterogeneity promotes robust learning.
\emph{Nature Communications}, 12:5791, 2021.
\bibitem{amir2017dvs} A. Amir et al. A low power, fully event-based gesture
recognition system. \emph{CVPR}, 2017.
\bibitem{neftci2019surrogate} E. O. Neftci, H. Mostafa, F. Zenke. Surrogate
gradient learning in spiking neural networks. \emph{IEEE Signal Processing
Magazine}, 36(6), 2019.
\bibitem{zenke2021remarkable} F. Zenke, T. P. Vogels. The remarkable
robustness of surrogate gradient learning for instilling complex function in
spiking neural networks. \emph{Neural Computation}, 33(4), 2021.
\bibitem{fang2021incorporating} W. Fang et al. Incorporating learnable
membrane time constant to enhance learning of spiking neural networks.
\emph{ICCV}, 2021.
\end{thebibliography}

\end{document}
